\section{床沈み込みが歩行制御に与える影響}

床沈み込みは，二足歩行ロボットの歩行制御に対して複数の影響を与える．第一に，支持脚の足先位置が計画された位置から鉛直方向にずれることで，運動学的拘束条件が変化する点が挙げられる．この影響は，関節角度や姿勢制御に影響を及ぼし，制御系が想定していない誤差を生じさせる．

第二に，床反力の立ち上がり特性が変化する点である．剛体床では接地と同時に十分な床反力が得られるのに対し，軟弱地面では床変形を伴うため，床反力の立ち上がりが遅れ，一時的に支持力が不足する可能性がある．この影響は特に片脚支持期において顕著であり，歩行安定性の低下につながる．